% Generated draft: requires mathematical and source review.
\documentclass[11pt]{article}
\usepackage[margin=25mm]{geometry}
\usepackage{fontspec}
\setmainfont{Latin Modern Roman}
\usepackage{amsmath,amssymb,mathtools,unicode-math}
\setmathfont{Latin Modern Math}
\usepackage{enumitem,needspace,xurl,hyperref}
\setlength{\parindent}{0pt}
\setlength{\parskip}{5pt}
\setlength{\emergencystretch}{3em}
\newcommand{\N}{\mathbb{N}}
\newcommand{\Z}{\mathbb{Z}}
\newcommand{\Q}{\mathbb{Q}}
\newcommand{\R}{\mathbb{R}}
\newcommand{\C}{\mathbb{C}}
\newcommand{\F}{\mathbb{F}}
\newcommand{\A}{\mathbb{A}}
\newcommand{\PP}{\mathbb{P}}
\newcommand{\E}{\mathbb{E}}
\newcommand{\eps}{\varepsilon}
\newcommand{\dd}{\,\mathrm{d}}
\begin{document}
\section*{UnsolvedMath ID: 11000304}
\textbf{Source problem code:} catalogue entry
\subsection{Definitions and mathematical statement}
For a closed oriented genus-\(g\) surface, let \(I_g\) be the Torelli kernel of the mapping-class action on \(H_1(\Sigma_g;\Z)\). If \(p:E\to B\operatorname{Mod}_g\) is the universal surface bundle and \(e\) its vertical Euler class, the usual Mumford--Morita--Miller class is \(e_j=p_!(e^{j+1})\in H^{2j}(B\operatorname{Mod}_g;\Q)\). Use the same notation for its restriction to \(BI_g\).

Prove or disprove nontriviality of the even classes
\[
 e_{2i}\in H^{4i}(I_g;\Q)
\]
in a suitable stable genus range. The standard large-genus reading is: for each positive \(i\), does there exist \(g_0(i)\) such that the restricted class is nonzero for every \(g\geq g_0(i)\)? The source supplies no explicit range and does not assert homological stability of the Torelli groups themselves. Nonvanishing on the full mapping class group alone does not settle its restriction.

The source attributes to Johnson [40] the large rational abelianization \(H_1(I_g;\Q)\cong U_\Q\) for \(g\geq3\), with \(U_\Q\) defined earlier but absent here. It contrasts this with the counterfactual perfect-kernel situation, where it suggests applying Quillen plus constructions to the missing extension (1) and using a homotopy exact sequence. This is not a valid substitute for treating the actual nonperfect kernel.

It cites Igusa [38, Corollary 8.5.17] for a close connection with nontriviality of higher Franz--Reidemeister torsion in
\(H^{4i}(I\operatorname{Out}_n;\R)\), where
\(I\operatorname{Out}_n=\ker(\operatorname{Out}(F_n)\to\operatorname{GL}(n,\Z))\), denoted \(\operatorname{Out}^{h}F_n\) by Igusa. The precise comparison and Problem 4.4 of Section 4 are missing. Related Sakasai [95] and foundational Hain [29] work are mentioned; the announced following Torelli-structure problems are not supplied.
\subsection{Short English statement}
Determine whether the even MMM classes remain nonzero on the Torelli group in large genus, retaining the nonperfect-kernel obstacle and the reported higher-torsion connection.
\subsection{Sources}
\begin{itemize}
\item ``Cohomological structure of the mapping class group and beyond,'' in \emph{Problems on Mapping Class Groups and Related Topics}, edited by Farb, Problem 3.1 and continuation, supplied PDF page 361. \url{https://www.math.uchicago.edu/~farb/papers/mcgbook.pdf}. Recorded access 2026-07-29; extraction pdf-font-structure.
\item Johnson [40]; Igusa [38], particularly Corollary 8.5.17 and the source's Section 4, Problem 4.4; Sakasai [95]; Hain [29]. Full bibliography entries, extension (1) and the referenced problem are absent.
\item Local UnsolvedMath registry, \texttt{data/unsolvedmath/problems.json}, permanent ID \(11000304\), exact source code \texttt{catalogue entry}; supplied OPEN-TRIAGE summary distinguishes full-group MMM nontriviality from the Torelli restriction.
\end{itemize}
\end{document}
